\section{Research Notes and Selected Sources}

\subsection{Theoretical Grounding and Practitioner Tools}

This chapter synthesizes empirical scholarship from organizational behavior, strategic management, corporate governance, and entrepreneurial economics with practitioner decision frameworks. The ambition vector, mission-discovery interview, boundary register, business-model constitution, idea cartography, and 14-day architecture sprint are practical heuristics designed to make strategic trade-offs explicit and defend the founder against premature operational certainty.

\subsection{Primary Sources and Scholarly References}

\begin{enumerate}[leftmargin=12pt,itemsep=3pt,topsep=1pt,parsep=0pt]
  \item \textbf{Collins, J. C., \& Porras, J. I. (1996).} ``Building Your Company's Vision.'' \textit{Harvard Business Review}, September--October 1996. Establishes the vital conceptual separation between enduring core purpose and foundational values, and the constantly adapting operational strategies and tactical practices of an enterprise.
  \item \textbf{Porter, M. E. (1996).} ``What Is Strategy?'' \textit{Harvard Business Review}, November--December 1996. The foundational economic treatise on strategic positioning: demonstrating that operational effectiveness is not strategy, and that strategy requires clear trade-offs, mutually reinforcing activities, and the discipline of saying no.
  \item \textbf{Ebrahim, A., Battilana, J., \& Mair, J. (2014).} ``The governance of social enterprises: Mission drift and accountability challenges in hybrid organizations.'' \textit{Research in Organizational Behavior}, 34, 81--100. Rigorous academic analysis of structural governance tensions in organizations combining commercial market revenue with social and educational missions.
  \item \textbf{Rumelt, R. P. (2011).} \textit{Good Strategy/Bad Strategy: The Difference and Why It Matters}. Crown Business. Introduces the classic strategy kernel: diagnosis of the central obstacle, guiding policy, and coherent coordinated action.
  \item \textbf{Sarasvathy, S. D. (2001).} ``Causation and Effectuation: Toward a Theoretical Shift from Economic Inevitability to Entrepreneurial Contingency.'' \textit{Academy of Management Review}, 26(2), 243--263. The foundational paper on effectuation, beginning with immediate available means, strictly bounding downside exposure via affordable loss, and leveraging stakeholder commitments.
  \item \textbf{March, J. G. (1991).} ``Exploration and Exploitation in Organizational Learning.'' \textit{Organization Science}, 2(1), 71--87. The definitive model examining the precarious balance between exploiting existing organizational efficiencies and exploring uncertain future possibilities.
  \item \textbf{University Grants Commission (India).} \textit{Official Statutory Regulations and FAQs on Degree-Awarding Institutions}. Government of India. Details the mandatory statutory frameworks governing university establishment, deemed-to-be-university status under Section~3 of the UGC Act, and the exclusive legal right to confer recognized academic degrees.
  \item \textbf{Hamel, G., \& Prahalad, C. K. (1989).} ``Strategic Intent.'' \textit{Harvard Business Review}, May--June 1989. Analyzes the motivating power of ambitious, multi-decade organizational intent balanced against feasible resource deployment and downside risk.
  \item \textbf{Battilana, J., \& Lee, M. (2014).} ``Advancing Research on Hybrid Organizing: Insights from the Study of Social Enterprises.'' \textit{Academy of Management Annals}, 8(1), 397--441. Comprehensive synthesis of the institutional logics and governance architectures required to prevent mission drift in hybrid organizations.
  \item \textbf{Graham, P. (2012).} ``How to Get Startup Ideas.'' Available at \texttt{paulgraham.com}. Emphasizes solving authentic personal problems, developing organic domain insight, and serving real users rather than manufacturing contrived commercial concepts.
\end{enumerate}

\subsection{Comparative Institutional Models}

For founders examining the structural evolution of private, multidisciplinary research institutions, the public governance structures and founding histories of \textbf{Stanford University}, \textbf{Plaksha University}, and \textbf{Ashoka University} provide instructive case studies.

Each demonstrates a distinctive configuration of philanthropic endowments, corporate partnership models, interdisciplinary STEM pedagogy, and long-term academic governance. They must be studied as evolutionary architectures, however, rather than blueprints to be casually copied by an early-stage venture without comparable capital.

\subsection{Epistemological Limits of This Chapter}

This chapter does \textbf{not} establish the definitive flagship product for Bayes Institute, its initial tuition pricing power, verified student learning outcomes, corporate willingness to pay, formal legal entity structure, or a specific calendar date for campus land acquisition.

Those operational questions require empirical customer discovery, pilot delivery data, financial modeling, and formal legal counsel.

Its narrower, essential achievement has been to construct the founder's decision-making constitution: \textbf{defining what the institution exists to achieve, which commercial practices are categorically forbidden, which operational options remain open, and what empirical evidence must guide the opening move.}
